% Ad Familiares 5.12 (ad-familiares-05-12)
% Source: https://raw.githubusercontent.com/PerseusDL/canonical-latinLit/master/data/phi0474/phi056/phi0474.phi056.perseus-lat1.xml
% Fetched by scripts/fetch_latin.py

% Dateline (Perseus): Scr. Ante ut videtur, m. Iun. a. 698 (56) .

\ciceroLetterOpener{M. CICERO S. D. L. LVCCEIO Q. F. .}

\ciceroSection{1} Coram me tecum eadem haec agere saepe conantem deterruit pudor quidam paene subrusticus, quae nunc expromam absens audacius; epistula enim non erubescit. ardeo cupiditate incredibili neque, ut ego arbitror, reprehendenda, nomen ut nostrum scriptis inlustretur et celebretur tuis. quod etsi mihi saepe ostendisti te esse facturum, tamen ignoscas velim huic festinationi meae. genus enim scriptorum tuorum etsi erat semper a me vehementer exspectatum, tamen vicit opinionem meam meque ita vel cepit vel incendit, ut cuperem quam celerrime res nostras monumentis commendari tuis. neque enim me solum commemoratio posteritatis ad spem quandam immortalitatis rapit, sed etiam illa cupiditas, ut vel auctoritate testimoni tui vel indicio benevolentiae vel suavitate ingeni vivi perfruamur.

\ciceroSection{2} neque tamen, haec cum scribebam, eram nescius quantis oneribus premerere susceptarum rerum et iam institutarum sed, quia videbam Italici belli et civilis historiam iam a te paene esse perfectam, dixeras autem mihi te reliquas res ordiri, desse mihi nolui quin te admonerem, ut cogitares, coniunctene malles cum reliquis rebus nostra contexere an, ut multi Graeci fecerunt, Callisthenes Phocicum bellum, Timaeus Pyrrhi, Polybius Numantinum, qui omnes a perpetuis suis historiis ea, quae dixi, bella separaverunt, tu quoque item civilem coniurationem ab hostilibus externisque bellis seiungeres. equidem ad nostram laudem non multum video interesse, sed ad properationem meam quiddam interest non te exspectare, dum ad locum venias, ac statim causam illam totam et tempus arripere; et simul, si uno in argumento unaque in persona mens tua tota versabitur, cerno iam animo, quanto omnia uberiora atque ornatiora futura sint. neque tamen ignoro quam impudenter faciam, qui primum tibi tantum oneris imponam (potest enim mihi denegare occupatio tua), deinde etiam ut ornes me postulem. quid si illa tibi non tanto opere videntur ornanda?

\ciceroSection{3} sed tamen, qui semei verecundiae finis transierit, eum bene et naviter oportet esse impudentem. itaque te plane etiam atque etiam rogo, ut et ornes ea vehementius etiam quam fortasse sentis, et in eo leges historiae neglegas gratiamque illam, de qua suavissime quodam in prooemio scripsisti, a qua te deflecti non magis potuisse demonstras quam Herculem Xenophontium illum a voluptate, eam, si me tibi vehementius commendabit, ne aspernere amorique nostro plusculum etiam, quam concedet veritas, largiare. quod si te adducemus ut hoc suscipias, erit, ut mihi persuadeo, materies digna facultate et copia tua.

\ciceroSection{4} A principio enim coniurationis usque ad reditum nostrum videtur mihi modicum quoddam corpus confici posse, in quo et illa poteris uti civilium commutationum scientia vel in explicandis causis rerum novarum vel in remediis incommodorum, cum et reprehendes ea, quae vituperanda duces, et quae placebunt exponendis rationibus comprobabis et, si liberius, ut consuesti, agendum putabis, multorum in nos perfidiam, insidias, proditionem notabis. multam etiam casus nostri varietatem tibi in scribendo suppeditabunt plenam cuiusdam voluptatis, quae vehementer animos hominum in legendo te scriptore tenere possit. nihil est enim aptius ad delectationem lectoris quam temporum varietates fortunaeque vicissitudines. quae etsi nobis optabiles in experiendo non fuerunt, in legendo tamen erunt iucundae; habet enim praeteriti doloris secura recordatio delectationem;

\ciceroSection{5} ceteris vero nulla perfunctis propria molestia, casus autem alienos sine ullo dolore intuentibus etiam ipsa misericordia est iucunda. quem enim nostrum ille moriens apud Mantineam Epaminondas non cum quadam miseratione delectat? qui tum denique sibi evelli iubet spiculum, postea quam ei percontanti dictum est clipeum esse salvum, ut etiam in vulneris dolore aequo animo cum laude moreretur. cuius studium in legendo non erectum Themistocli fuga redituque retinetur? etenim ordo ipse annalium mediocriter nos retinet quasi enumeratione fastorum at viri saepe excellentis ancipites variique casus habent admirationem, exspectationem, laetitiam, molestiam, spem, timorem; si vero exitu notabili concluduntur, expletur animus iucundissima lectionis voluptate.

\ciceroSection{6} quo mihi acciderit optatius, si in hac sententia fueris ut a continentibus tuis scriptis, in quibus perpetuam rerum gestarum historiam complecteris, secernas hanc quasi fabulam rerum eventorumque nostrorum. habet enim varios actus mutationesque et consiliorum et temporum. ac non vereor ne adsentatiuncula quadam aucupari tuam gratiam videar, quom hoc demonstrem, me a te potissimum ornari celebrarique velle. neque enim tu is es qui, quid sis, nescias et qui non eos magis, qui te non admirentur, invidos quam eos, qui laudent, adsentatores arbitrere; neque autem ego sum ita demens, ut me sempiternae gloriae per eum commendari velim, qui non ipse quoque in me commendando propriam ingeni gloriam consequatur.

\ciceroSection{7} neque enim Alexander ille gratiae causa ab Apelle potissimum pingi et a Lysippo fingi volebat, sed quod illorum artem cum ipsis tum etiam sibi gloriae fore putabat. atque illi artifices corporis simulacra ignotis nota faciebant; quae vel si nulla sint, nihilo sint tamen obscuriores clari viri. nec minus est Spartiates Agesilaus ille perhibendus, qui neque pictam neque fictam imaginem suam passus est esse, quam qui in eo genere laborarunt; unus enim Xenophontis libellus in eo rege laudando facile omnis imagines omnium statuasque superavit. atque hoc praestantius mihi fuerit et ad laetitiam animi et ad memoriae dignitatem, si in tua scripta pervenero quam si in ceterorum, quod non ingenium mihi solum suppeditatum fuerit tuum, sicut Timoleonti a Timaeo aut ab Herodoto Themistocli, sed etiam auctoritas clarissimi et spectatissimi viri et in rei p. maximis gravissimisque causis cogniti atque in primis probati, ut mihi non solum praeconium, quod, cum in Sigeum venisset, Alexander ab Homero Achilli tributum esse dixit, sed etiam grave testimonium impertitum clari hominis magnique videatur. placet enim Hector ille mihi Naevianus, qui non tantum 'laudari' se laetatur, sed addit etiam 'a laudato viro.'

\ciceroSection{8} quod si a te non impetro, hoc est, si quae te res impedierit (neque enim fas esse arbitror quicquam me rogantem abs te non impetrare), cogar fortasse facere, quod non nulli saepe reprehendunt, scribam ipse de me, multorum tamen exemplo et clarorum virorum. sed , quod te non fugit, haec sunt in hoc genere vitia: et verecundius ipsi de sese scribant necesse est, si quid est laudandum, et praetereant, si quid reprehendendum est. accedit etiam ut minor sit fides, minor auctoritas, multi denique reprehendant et dicant verecundiores esse praecones ludorum gymnicorum, qui cum ceteris coronas imposuerint victoribus eorumque nomina magna voce pronuntiarint, cum ipsi ante ludorum missionem corona donentur, alium praeconem adhibeant, ne sua voce se ipsi victores esse praedicent.

\ciceroSection{9} haec nos vitare cupimus et, si recipis causam nostram, vitabimus, idque ut facias rogamus. ac ne forte mirere cur, cum mihi saepe ostenderis te accuratissime nostrorum temporum consilia atque eventus litteris mandaturum, a te id nunc tanto opere et tam multis verbis petamus, illa nos cupiditas incendit, de qua initio scripsi, festinationis, quod alacres animo sumus, ut et ceteri viventibus nobis ex libris tuis nos cognoscant et nosmet ipsi vivi gloriola nostra perfruamur.

\ciceroSection{10} his de rebus quid acturus sis, si tibi non est molestum, rescribas mihi velim. si enim suscipis causam, conficiam commentarios rerum omnium, sin autem differs me in tempus aliud, coram tecum loquar. tu interea non cessabis et ea, quae habes instituta, perpolies nosque diliges.
